\documentclass[a4paper,12pt,twoside]{memoir}

% Castellano
\usepackage[spanish,es-tabla]{babel}
\selectlanguage{spanish}
\usepackage[utf8]{inputenc}
\usepackage[T1]{fontenc}
\usepackage{lmodern} % Scalable font
\usepackage{microtype}
\usepackage{placeins}

\RequirePackage{booktabs}
\RequirePackage[table]{xcolor}
\RequirePackage{xtab}
\RequirePackage{multirow}

% Links
\PassOptionsToPackage{hyphens}{url}\usepackage[colorlinks]{hyperref}
\hypersetup{
	allcolors = {red}
}

% Ecuaciones
\usepackage{amsmath}

% Rutas de fichero / paquete
\newcommand{\ruta}[1]{{\sffamily #1}}

% Párrafos
\nonzeroparskip

% Huérfanas y viudas
\widowpenalty100000
\clubpenalty100000

% Imágenes

% Comando para insertar una imagen en un lugar concreto.
% Los parámetros son:
% 1 --> Ruta absoluta/relativa de la figura
% 2 --> Texto a pie de figura
% 3 --> Tamaño en tanto por uno relativo al ancho de página
\usepackage{graphicx}
\newcommand{\imagen}[3]{
	\begin{figure}[!h]
		\centering
		\includegraphics[width=#3\textwidth]{#1}
		\caption{#2}\label{fig:#1}
	\end{figure}
	\FloatBarrier
}

% Comando para insertar una imagen sin posición.
% Los parámetros son:
% 1 --> Ruta absoluta/relativa de la figura
% 2 --> Texto a pie de figura
% 3 --> Tamaño en tanto por uno relativo al ancho de página
\newcommand{\imagenflotante}[3]{
	\begin{figure}
		\centering
		\includegraphics[width=#3\textwidth]{#1}
		\caption{#2}\label{fig:#1}
	\end{figure}
}

% El comando \figura nos permite insertar figuras comodamente, y utilizando
% siempre el mismo formato. Los parametros son:
% 1 --> Porcentaje del ancho de página que ocupará la figura (de 0 a 1)
% 2 --> Fichero de la imagen
% 3 --> Texto a pie de imagen
% 4 --> Etiqueta (label) para referencias
% 5 --> Opciones que queramos pasarle al \includegraphics
% 6 --> Opciones de posicionamiento a pasarle a \begin{figure}
\newcommand{\figuraConPosicion}[6]{%
  \setlength{\anchoFloat}{#1\textwidth}%
  \addtolength{\anchoFloat}{-4\fboxsep}%
  \setlength{\anchoFigura}{\anchoFloat}%
  \begin{figure}[#6]
    \begin{center}%
      \Ovalbox{%
        \begin{minipage}{\anchoFloat}%
          \begin{center}%
            \includegraphics[width=\anchoFigura,#5]{#2}%
            \caption{#3}%
            \label{#4}%
          \end{center}%
        \end{minipage}
      }%
    \end{center}%
  \end{figure}%
}

%
% Comando para incluir imágenes en formato apaisado (sin marco).
\newcommand{\figuraApaisadaSinMarco}[5]{%
  \begin{figure}%
    \begin{center}%
    \includegraphics[angle=90,height=#1\textheight,#5]{#2}%
    \caption{#3}%
    \label{#4}%
    \end{center}%
  \end{figure}%
}
% Para las tablas
\newcommand{\otoprule}{\midrule [\heavyrulewidth]}
%
% Nuevo comando para tablas pequeñas (menos de una página).
\newcommand{\tablaSmall}[5]{%
 \begin{table}
  \begin{center}
   \rowcolors {2}{gray!35}{}
   \begin{tabular}{#2}
    \toprule
    #4
    \otoprule
    #5
    \bottomrule
   \end{tabular}
   \caption{#1}
   \label{tabla:#3}
  \end{center}
 \end{table}
}

%
% Nuevo comando para tablas pequeñas (menos de una página).
\newcommand{\tablaSmallSinColores}[5]{%
 \begin{table}[H]
  \begin{center}
   \begin{tabular}{#2}
    \toprule
    #4
    \otoprule
    #5
    \bottomrule
   \end{tabular}
   \caption{#1}
   \label{tabla:#3}
  \end{center}
 \end{table}
}

\newcommand{\tablaApaisadaSmall}[5]{%
\begin{landscape}
  \begin{table}
   \begin{center}
    \rowcolors {2}{gray!35}{}
    \begin{tabular}{#2}
     \toprule
     #4
     \otoprule
     #5
     \bottomrule
    \end{tabular}
    \caption{#1}
    \label{tabla:#3}
   \end{center}
  \end{table}
\end{landscape}
}

%
% Nuevo comando para tablas grandes con cabecera y filas alternas coloreadas en gris.
\newcommand{\tabla}[6]{%
  \begin{center}
    \tablefirsthead{
      \toprule
      #5
      \otoprule
    }
    \tablehead{
      \multicolumn{#3}{l}{\small\sl continúa desde la página anterior}\\
      \toprule
      #5
      \otoprule
    }
    \tabletail{
      \hline
      \multicolumn{#3}{r}{\small\sl continúa en la página siguiente}\\
    }
    \tablelasttail{
      \hline
    }
    \bottomcaption{#1}
    \rowcolors {2}{gray!35}{}
    \begin{xtabular}{#2}
      #6
      \bottomrule
    \end{xtabular}
    \label{tabla:#4}
  \end{center}
}

%
% Nuevo comando para tablas grandes con cabecera.
\newcommand{\tablaSinColores}[6]{%
  \begin{center}
    \tablefirsthead{
      \toprule
      #5
      \otoprule
    }
    \tablehead{
      \multicolumn{#3}{l}{\small\sl continúa desde la página anterior}\\
      \toprule
      #5
      \otoprule
    }
    \tabletail{
      \hline
      \multicolumn{#3}{r}{\small\sl continúa en la página siguiente}\\
    }
    \tablelasttail{
      \hline
    }
    \bottomcaption{#1}
    \begin{xtabular}{#2}
      #6
      \bottomrule
    \end{xtabular}
    \label{tabla:#4}
  \end{center}
}

%
% Nuevo comando para tablas grandes sin cabecera.
\newcommand{\tablaSinCabecera}[5]{%
  \begin{center}
    \tablefirsthead{
      \toprule
    }
    \tablehead{
      \multicolumn{#3}{l}{\small\sl continúa desde la página anterior}\\
      \hline
    }
    \tabletail{
      \hline
      \multicolumn{#3}{r}{\small\sl continúa en la página siguiente}\\
    }
    \tablelasttail{
      \hline
    }
    \bottomcaption{#1}
  \begin{xtabular}{#2}
    #5
   \bottomrule
  \end{xtabular}
  \label{tabla:#4}
  \end{center}
}



\definecolor{cgoLight}{HTML}{EEEEEE}
\definecolor{cgoExtralight}{HTML}{FFFFFF}

%
% Nuevo comando para tablas grandes sin cabecera.
\newcommand{\tablaSinCabeceraConBandas}[5]{%
  \begin{center}
    \tablefirsthead{
      \toprule
    }
    \tablehead{
      \multicolumn{#3}{l}{\small\sl continúa desde la página anterior}\\
      \hline
    }
    \tabletail{
      \hline
      \multicolumn{#3}{r}{\small\sl continúa en la página siguiente}\\
    }
    \tablelasttail{
      \hline
    }
    \bottomcaption{#1}
    \rowcolors[]{1}{cgoExtralight}{cgoLight}

  \begin{xtabular}{#2}
    #5
   \bottomrule
  \end{xtabular}
  \label{tabla:#4}
  \end{center}
}



\graphicspath{ {./img/} }

% Capítulos
\chapterstyle{bianchi}
\newcommand{\capitulo}[2]{
	\setcounter{chapter}{#1}
	\setcounter{section}{0}
	\setcounter{figure}{0}
	\setcounter{table}{0}
	\chapter*{\thechapter.\enskip #2}
	\addcontentsline{toc}{chapter}{\thechapter.\enskip #2}
	\markboth{#2}{#2}
}

% Apéndices
\renewcommand{\appendixname}{Apéndice}
\renewcommand*\cftappendixname{\appendixname}

\newcommand{\apendice}[1]{
	%\renewcommand{\thechapter}{A}
	\chapter{#1}
}

\renewcommand*\cftappendixname{\appendixname\ }

% Formato de portada
\makeatletter
\usepackage{xcolor}
\newcommand{\tutor}[1]{\def\@tutor{#1}}
\newcommand{\course}[1]{\def\@course{#1}}
\definecolor{cpardoBox}{HTML}{E6E6FF}
\def\maketitle{
  \null
  \thispagestyle{empty}
  % Cabecera ----------------
\noindent\includegraphics[width=\textwidth]{cabecera}\vspace{1cm}%
  \vfill
  % Título proyecto y escudo informática ----------------
  \colorbox{cpardoBox}{%
    \begin{minipage}{.8\textwidth}
      \vspace{.5cm}\Large
      \begin{center}
      \textbf{TFG del Grado en Ingeniería Informática}\vspace{.6cm}\\
      \textbf{\LARGE\@title{}}
      \end{center}
      \vspace{.2cm}
    \end{minipage}

  }%
  \hfill\begin{minipage}{.20\textwidth}
    \includegraphics[width=\textwidth]{escudoInfor}
  \end{minipage}
  \vfill
  % Datos de alumno, curso y tutores ------------------
  \begin{center}%
  {%
    \noindent\LARGE
    Presentado por \@author{}\\ 
    en Universidad de Burgos --- \@date{}\\
    Tutor: \@tutor{}\\
  }%
  \end{center}%
  \null
  \cleardoublepage
  }
\makeatother

\newcommand{\nombre}{Víctor De Marco Velasco}
\newcommand{\titulo}{Estrategias de seguridad para dispositivos Internet de las cosas industriales (IioT)}
%%% cambio de comando

% Datos de portada
\title{\titulo}
\author{\nombre}
\tutor{Carlos Cambra Baseca}
\date{\today}

\begin{document}

\maketitle


\newpage\null\thispagestyle{empty}\newpage


%%%%%%%%%%%%%%%%%%%%%%%%%%%%%%%%%%%%%%%%%%%%%%%%%%%%%%%%%%%%%%%%%%%%%%%%%%%%%%%%%%%%%%%%
\thispagestyle{empty}


\noindent\includegraphics[width=\textwidth]{cabecera}\vspace{1cm}

\noindent D. nombre tutor, profesor del departamento de nombre departamento, área de nombre área.

\noindent Expone:

\noindent Que el alumno D. \nombre, con DNI 71706956V, ha realizado el Trabajo final de Grado en Ingeniería Informática titulado \titulo. 

\noindent Y que dicho trabajo ha sido realizado por el alumno bajo la dirección del que suscribe, en virtud de lo cual se autoriza su presentación y defensa.

\begin{center} %\large
En Burgos, {\large \today}
\end{center}

\vfill\vfill\vfill

% Author and supervisor
\begin{minipage}{0.45\textwidth}
\begin{flushleft} %\large
Vº. Bº. del Tutor:\\[2cm]
D. nombre tutor
\end{flushleft}
\end{minipage}
\hfill
\begin{minipage}{0.45\textwidth}
\begin{flushleft} %\large
Vº. Bº. del co-tutor:\\[2cm]
D. nombre co-tutor
\end{flushleft}
\end{minipage}
\hfill

\vfill

% para casos con solo un tutor comentar lo anterior
% y descomentar lo siguiente
%Vº. Bº. del Tutor:\\[2cm]
%D. nombre tutor


\newpage\null\thispagestyle{empty}\newpage




\frontmatter

% Abstract en castellano
\renewcommand*\abstractname{Resumen}
\begin{abstract}
Este proyecto propone el desarrollo de diferentes aplicaciones web con el propósito de mejorar la seguridad de los dispositivos IoT y la interacción del usuario con los datos recopilados por dichos dispositivos. Una página web se encarga de recibir los datos, comprobar si dichos datos han sido realmente enviados por nuestro dispositivo o si es posible que hayan sido generados por un tercero con intenciones maliciosas y, una vez comprobado, proceda a almacenarlos. La segunda página web está más enfocada al usuario, permitiéndole interactuar con los datos recopilados e incluso almacenar/analizar los suyos propios.
\end{abstract}

\renewcommand*\abstractname{Descriptores}
\begin{abstract}
internet de las cosas, LoRaWAN, Reglas de asociación, The things network, Detección de anomalías,Python, Flujos de red\ldots
\end{abstract}

\clearpage

% Abstract en inglés
\renewcommand*\abstractname{Abstract}
\begin{abstract}
This project proposes the development of different web apps with the aim of improving the security of IoT devices and the user's interaction with the data collected by these devices. One web page is responsible for receiving the data, verifying whether the data was genuinely sent by our device or if it could have been generated by a third party with malicious intent, and, once verified, storing it. The second web page is more user-focused, allowing the user to interact with the collected data and even store and analyze their own data..
\end{abstract}

\renewcommand*\abstractname{Keywords}
\begin{abstract}
IoT, LoRaWAN, Association rules, The things network, Anomaly detection,Python, Network flows\ldots
\end{abstract}

\clearpage

% Indices
\tableofcontents

\clearpage

\listoffigures

\clearpage

\listoftables
\clearpage

\mainmatter
\capitulo{1}{Introducción}

Descripción del contenido del trabajo y del estructura de la memoria y del resto de materiales entregados.

\capitulo{2}{Objetivos del proyecto}

Este apartado explica de forma precisa y concisa cuales son los objetivos que se persiguen con la realización del proyecto. Se puede distinguir entre los objetivos marcados por los requisitos del software a construir, los objetivos de carácter técnico que plantea a la hora de llevar a la práctica el proyecto y los objetivos personales del autor.

\section {Objetivos Generales}
\begin{itemize}
   \item Diseñar e implementar una solución completa que permite recoger, analizar, almacenar y visualizar los datos generados por un dispositivo IoT conectado a una red LoRaWAN.
    \item Proporcionar a usuarios no técnicos una interfaz web intuitiva desde la que puedan interactuar con los datos, analizar su contenido y detectar posibles anomalías.
\end{itemize}
\section {Objetivos Técnicos}
\begin{itemize}
   \item Configurar y poner en funcionamiento una infraestructura IoT basada en TTN (The Things Network), incluyendo un gateway LoRaWAN Dragino y un dispositivo MerryIoT.
    \item Desarrollar una aplicación web de recepción que utilice un webhook de TTN y un túnel de Cloudflare para recolectar en tiempo real los datos generados por el dispositivo.
    \item Implementar un sistema de evaluación heurística de los paquetes recibidos para clasificarlos automáticamente  como reales o infectados según sus características.
    \item Almacenar los datos evaluados en un archivo CSV, creando así un dataset que será utilizado a futuro.
    \item Desarrollar una segunda aplicación web de gestión y visualización, donde el usuario pueda:
    \begin{itemize}
     \item Subir y analizar sus propios archivos CSV.
     \item Visualizar gráficamente información clave como temperatura o humedad.
     \item Simular el envío de nuevos paquetes a la aplicación web de recepción.
     \end{itemize}
\end{itemize}
\section {Objetivos Personales} 
\begin{itemize}
   \item Realizar el desarrollo completo de un proyecto con sus diferentes apartados, mejorando así mis conocimientos sobre la gestión de proyectos.
    \item Aumentar mis conocimientos sobre los dispositivos IoT y el ecosistema que los rodea (TTN , LoRaWAN, ...)
    \item Aprender a desarrollar webs en Python mediante frameworks como Flask.
    \item Familiarizarme con el uso de software de control de versiones como Git y con Github la plataforma utilizada para alojar repositorios.
\end{itemize}
\capitulo{3}{Conceptos teóricos}

Para comprender correctamente las herramientas y tecnicas utilizas para cumplir los objetivos mencionados anteriormente es necesario comprender teóricamente estos conceptos, las redes de dispositivos IoT, la transmisión de datos mediante protocolos LoRaWAN y la recogida y análisis de datos. Dichos conceptos pasaremos a explicarlos detalladamente a continuación:



\section{Internet de las Cosas (IoT)}
\subsection{Definición}
El \textbf{internet de las cosas (IoT)}~\cite{aws:iot} hace referencia a una red de dispositivos físicos interconectados capaces de recopilar, intercambiar y procesar datos a través de Internet sin intervención humana directa. Esta tecnología permite que objetos cotidianos como sensores de movimiento, termostatos, electrodomésticos, vehículos o sistemas industriales puedan responder al usuario de forma inteligente.

El concepto de IoT surge de la evolución de tecnologías como los chips de bajo consumo, las etiquetas RFID y las telecomunicaciones de alta velocidad, que han facilitado la miniaturización y el abaratamiento de los componentes electrónicos. Gracias a estos avances, actualmente es viable integrar esos sensores en los objetos de uso cotidiano.
\subsection{Cómo funciona}
Un sistema común de IoT funciona mediante la recopilación y el intercambio de datos en tiempo real. 

Un sistema del IoT tiene tres componentes:
\begin{itemize}
   \item \textbf{Dispositivo inteligente}: Se trata de un dispositivo, en nuestro caso el \textbf{MerryIoT MS10 Motion Detection}~\cite{MerryIoT:manual} , al que se le dotó de capacidades de computación. Recopila datos de su entorno o de las entradas de los usuarios y comunica los datos a través de Internet hacia la aplicación de IoT.
    \item \textbf{Aplicación de IoT}: Una aplicación de IoT es un conjunto de servicios y software que integra los datos recibidos de varios dispositivos de IoT. Nuestra aplicación utiliza esos datos recibidos para analizarlos y tomar una decisión referente así el paquete recibido puede o no a ver sido infectado\footnote{Se dice infectado de un paquete de datos recibido por la aplicación IoT, que tras ser analizado muestra características sospechosas que nos dan a entender que ese paquete puede ser fruto de un tercero con intenciones maliciosas}.
    \item \textbf{Una interfaz de usuario gráfica}: El dispositivo de IoT puede administrarse a través de una interfaz de usuario gráfica. En nuestro caso un sitio web que puede utilizarse para analizar y gestionar archivos de datos recopilados por estos dispositivos.
\end{itemize}

\imagen{Dia_IoT}{Diagrama de un sistema IoT (con soporte movil)~\cite{IoT:diaimg}}{1}


\subsection{IioT o IoT industrial}
El \textbf{Internet Industrial de las Cosas (IIoT)}~\cite{ib:iiot} es el conjunto de sensores, instrumentos y dispositivos autónomos conectados a través de Internet a aplicaciones industriales. Esta red permite recopilar datos, realizar análisis y optimizar la producción, aumentando la eficiencia y reduciendo los costes del proceso de fabricación y prestación de servicios. Las aplicaciones industriales son ecosistemas tecnológicos completos que conectan dispositivos y a estos con las personas que gestionan los procesos en líneas de montaje, logística o distribución a gran escala.

La diferencia entre el Internet de las Cosas (IoT) y su versión industrial (IIoT) es que mientras el IoT está enfocado a servicios para los consumidores, el IIoT se concentra en aumentar la seguridad y la eficiencia en los centros de producción. Aplicado a nuestro proyecto al medir factores como temperatura y humedad podría aplicarse en un entorno referente a la industria de la agricultura donde es relevante llevar un control de esos factores.



\section{Redes LoRa y protocolo LoRaWAN}

\subsection{LoRa}

\subsubsection{Definición}

\textbf{LoRa (Long Range)}~\cite{Moko:Lora} es una tecnología de modulación inalámbrica que destaca por su bajo consumo y gran alcance, haciéndola ideal para aplicaciones IoT. Sin embargo, LoRa por sí sola no especifica cómo gestionar la comunicación entre dispositivos.

\subsubsection{¿Cómo funciona?}
Una arquitectura de red de Lorawan consiste en nodos Lora, Puertas de enlace de Lora, el servidor de red y el servidor de aplicaciones.

\imagen{LoRa_arq}{Arquitectura de una red LoRaWAN~\cite{TTN:Lorawan}}{1}

\begin{itemize}
   \item \textbf{Nodos LoRa}: Los nodos finales son los elementos de la red Lora donde se realiza el control o la detección. Normalmente están a prueba de baterías y se encuentran remotamente. Los nodos finales envían datos a cada puerta de enlace en su vecindad y transmiten datos en periódico.
    \item \textbf{LoRa gateway}: La puerta de enlace recibe los datos de los nodos Lora End y luego los canaliza a un servidor de red. Una puerta de enlace de Lora generalmente consiste en un módulo de radio Lora, un microprocesador, y un medio de conectividad a Internet.
    \item \textbf{Servidor de red}: El servidor de red administra la red. Filma paquetes duplicados causados por múltiples puertas de enlace que reciben los mismos datos, realiza controles de seguridad, administra el tráfico y el enrutamiento de la puerta de enlace, controlar la tasa adaptativa, y reenvía mensajes al servidor de aplicaciones.
    \item \textbf{Servidor de aplicaciones}:El servidor de aplicaciones procesa datos del servidor de red, Analiza datos del sensor, admite funciones como visualización de estado y alertas en tiempo real, y opcionalmente puede enviar respuestas al nodo final.
\end{itemize}

\subsection{LoRaWAN}

\subsubsection{Definición}

LoRaWAN~\cite{Mono:Lorawan} es un protocolo de comunicación LPWAN diseñado para conectar dispositivos IoT con bajo consumo de energía y largo alcance, como hemos mencionado antes se base en la tecnología LoRa permitiendo la transmisión de datos a distancias de hasta 15 km en entornos rurales y 2-5 km en entornos urbanos.

\subsubsection{Seguridad}

La seguridad es un aspecto fundamental y al que le hemos dado mucha importancia durante el desarrollo de este proyecto por ello voy a nombrar las principales medidas de seguridad de LoRaWAN:

\begin{enumerate}
   \item \textbf{Cifrado de extremo a extremo}: Los datos se transmiten encriptados con AES-128, lo que garantiza su confidencialidad.
    \item \textbf{Autenticación de dispositivos}: Cada sensor cuenta con una clave única para evitar accesos no autorizados~\cite{TTN:Sec}.
    \item \textbf{Gestión de claves de seguridad}: Se utilizan claves de sesión dinámicas para reducir el riesgo de ataques.
\end{enumerate}

Sin embargo, como cualquier tecnología inalámbrica, LoRaWAN no es inmune a ataques. Es importante implementar buenas prácticas como el uso de firewalls, segmentación de redes y monitorización continuo para detectar anomalías. En el caso de este proyecto para combatir los posibles ataques se ha optado por la monitorización continua y detección de anomalías. 

\subsection{Diferencias entre LoRa y LoRaWAN}

LoRa~\cite{Moko:Lora} describe la capa física inferior. LoRaWAN es un protocolo que describe las capas superiores de la red. LoRaWAN es un control de acceso a medios basado en la nube (MAC\footnote{Control de Acceso al Medio es una subcapa de la capa de enlace de datos en el modelo OSI,que regula cómo los dispositivos acceden y comparten un medio de transmisión común}) protocolo de capa, pero actúa principalmente como un protocolo de capa de red para gestionar la comunicación entre dispositivos de nodo final y puertas de enlace LPWAN, como protocolo de dirección, mantenido por la Alianza LoRa.

LoRaWAN Define la arquitectura del sistema y el protocolo de comunicación para la red., mientras que la capa física LoRa permite el enlace de comunicación de largo alcance.

\imagen{Capas_LoRa}{Diagrama de capas LoRaWAN~\cite{Moko:Lora}}{1}

\subsection{Diferencias entre LoRa y otras tecnologías de transmisión inalámbricas}
LoRaWAN es un protocolo de comunicación idóneo para paquetes útiles de pequeño tamaño (como datos de sensores) a largas distancias. La modulación LoRa proporciona un alcance de comunicación significativamente mayor con anchos de banda bajos en comparación con otras tecnologías inalámbricas de transmisión de datos. La siguiente imagen muestra algunas tecnologías de acceso que pueden utilizarse para la transmisión inalámbrica de datos, comparando su ancho de banda con su rango de transmisión.

\imagen{dif_LoRa}{Gráfica comparativa de tecnologías de transmisión inalámbrica~\cite{TTN:Graf}}{1}


\section{Estructura de los paquetes y análisis heurístico)}

\subsection{Estructura de los paquetes}
Los paquetes enviados por el dispositivo MerryIoT, llegan codificados en un formato especificado en su manual de usuario~\cite{MerryIoT:manual}, para poder utilizar esos datos es necesario decodificarlos con un payload\footnote{ payload (o carga útil) es la parte del mensaje que contiene los datos de interés que se transmiten entre dos dispositivos, excluyendo los encabezados (headers), metadatos u otra información de control} formatter.

\imagen{payload}{Formato del payload del dispositivo utilizado en el proyecto~\cite{MerryIoT:manual}}{1}

Una vez decodificados, se puede comenzar a analizar su contenido y sacar conclusiones del mismo.

\subsection{Analisis del contenido}
En este proyecto se ha determinado que la mejor forma para clasificar los datos recibidos es mediante el uso de reglas heurísticas generadas a partir de reglas de asociación y patrones observados durante el desarrollo del proyecto.

\subsubsection{Reglas de asociación}
Las Reglas de Asociación~\cite{DTM:Rules} son un método utilizado en minería de datos para descubrir relaciones significativas entre variables dentro de grandes conjuntos de datos. Estas reglas permiten identificar patrones frecuentes o comportamientos recurrentes, generalmente en forma de implicaciones del tipo:

Si A ocurre, entonces B tiende a ocurrir.
(Ejemplo: Si un cliente compra pan, también compra jamón).

Para entender como se han obtenido estas reglas es importante aclarar antes unos conceptos:


\begin{itemize}
    \item \textbf{Soporte}: Número de instancias que la regla predice correctamente.
    
    soporte(A --> C) = soporte(A U C)
    
    \item \textbf{Confianza}: Cociente entre el soporte y el número de instancias para las cuales la reglas es aplicable (el antecedente es cierto.).
    
    confianza(A --> C) = soporte(A U C)/ soporte(A)
    
    \item \textbf{Item}: un par de atributos/valores.

    \item \textbf{Item Set}:  todos los items que aparecen en una regla.
    
\end{itemize}

Para obtener las reglas de asociación se busca superar un soporte y una confianza mínima.  

La ventaja de los algoritmos de reglas de asociación sobre los algoritmos más estándar de árboles de decisión es que las asociaciones pueden existir entre cualquiera de los atributos. Un algoritmo de árbol de decisión generará reglas con una única conclusión, mientras que los algoritmos de asociación tratan de buscar muchas reglas, cada una de las cuales puede tener una conclusión diferente~\cite{DTM:Aso}.

 A diferencia de un árbol de decisión, el conjunto de reglas de asociación no se puede usar directamente para realizar predicciones de mismo modo que puede hacerlo con un modelo estándar (como un árbol de decisión o una red neuronal). Esto se debe a las diversas conclusiones diferentes posibles de las reglas. Otro nivel de transformación es preciso para transformar las reglas de asociación en un conjunto de reglas de clasificación. Por tanto, las reglas de asociación producidas por algoritmos de asociación se conocen como modelos sin refinar. 

 Para refinar dicho modelo facilitado por las Reglas de asociación utilizamos patrones identificados durante la recopilación de los datos. 

\subsubsection{Reglas heurísticas}
Las reglas heurísticas son normas simples o criterios aproximados que se aplican para detectar comportamientos sospechosos o inusuales.Estas reglas se generan al analizar los patrones de comportamiento conocidos del dispositivo IoT.

El análisis heurístico es un proceso utilizado en campos como la seguridad informática (antivirus), por lo tanto me parecia bastante adecuado para este proyecto~\cite{Heu:Sec}.

Estas reglas heurísticas sumadas a las reglas de asociación conforman el conjunto encargado de analizar los datos recibidos y clasificarlos correctamente.













\capitulo{4}{Técnicas y herramientas}
Esta parte de la memoria tiene como objetivo presentar las técnicas metodológicas y las herramientas de desarrollo que se han utilizado para llevar a cabo el proyecto.

\section{Técnicas utilizadas}
\subsection{SCRUM}
SCRUM es un marco de trabajo ágil utilizado principalmente en el desarrollo de software, divide el trabajo en ciclos cortos y regulares llamados "sprints", normalmente de 1 a 4 semanas, en los que se desarrolla un incremento funcional del producto.

Para obtener mas información sobre la metodología Scrum utilizada durante el proyecto consultar el apartado A.2.

\section{Herramientas utilizadas}
\subsection{Hardware}
\subsubsection{Gateway Dragino LPS8N}
El LPS8N~\cite{Haw:Dra} es una gateway LoRaWAN de código abierto. Permite enlazar una red inalámbrica LoRa con una red IP mediante WiFi o Ethernet. La tecnología inalámbrica LoRa permite a los usuarios enviar datos y alcanzar distancias extremadamente largas con tasas de transmisión de datos bajas.

\imagen{Gateway_dragino}{Imagen del Gateway Dragino LPS8N~\cite{Haw:Dra}}{0.5}

\subsubsection{Sensor MerryIoT MS10 Motion Detection}
El MS10 es un sensor de movimiento que utiliza conectividad LoRaWAN para comunicar la presencia o ausencia de una persona, ademas de recopilar otros datos como temperatura y humedad.  Su uso previsto es colocar el sensor con una buena vista de una habitación para detectar si hay movimiento o no en el área.

El sensor está compuesto por un detector infrarrojo pasivo (PIR) y una lente de Fresnel. El cuerpo principal contiene la electrónica activa necesaria para detectar movimiento y transmitir cualquier cambio a través de una red LoRaWAN.

También cuenta con detectores de vibración e inclinación en caso de manipulación. Una vez que se detecta un evento, el sensor enviará un mensaje ascendente (uplink).

\imagen{MS10_device}{Imagen del dispositivo MS10~\cite{MerryIoT:manual}}{0.5}


\subsection{Servicios en la nube y conectividad}
\subsubsection{The Things Network (TTN)/The Things Stack (Sandbox)}
TTN es una infraestructura abierta para redes LoRaWAN. Permite registrar dispositivos, configurar gateways y gestionar los datos recibidos mediante webhooks, que envían automáticamente los paquetes recibidos hacia una URL determinada.

\imagen{TTN_logo}{Logotipo de The Things Network }{0.5}

The Things Stack~\cite{TTN:TTS} es un servidor LoRaWAN de nivel empresarial (que incluye tanto las funciones de Servidor de red como de Servidor de aplicaciones mencionadas en la arquitectura de referencia de LoRaWAN). Además, The Things Stack incluye servicios y herramientas para gestionar de forma segura millones de dispositivos LoRaWAN en entornos de producción.

En el proyecto se ha utilizado la variante The Things Stack Sandbox una versión gratuita y pública del servidor The Things Stack, ofrecida por The Things Network (TTN), diseñada para que desarrolladores, estudiantes e investigadores puedan probar, experimentar y aprender a implementar soluciones LoRaWAN sin necesidad de infraestructura propia

\imagen{TTS_logo}{Logotipo de The Things Stack Sandbox }{0.5}

\subsubsection{Cloudflare Tunnel}
Cloudflare Tunnel ~\cite{CDF:tun} es una herramienta que te proporciona una forma segura de conectar tus recursos a Cloudflare sin necesidad de una dirección IP pública enrutable. Con Tunnel, no envías tráfico a una IP externa; en su lugar, un servicio ligero (daemon) llamado cloudflared, que se ejecuta en tu infraestructura, establece conexiones salientes únicamente hacia la red global de Cloudflare.

En este proyecto se ha utilizado para exponer de forma segura tu servidor Flask local al exterior, permitiendo que The Things Network (TTN) pueda enviar datos directamente a tu webhook sin necesidad de:

\begin{enumerate}
    \item Configurar el router o abrir puertos manualmente.
    \item Tener una IP pública fija
    \item Contratar infraestructura externa (como un servidor en la nube)
\end{enumerate}

\imagen{CDF_logo}{Logotipo de Cloudflare }{0.5}

\subsection{Entorno de desarrollo}

\subsubsection{IntelliJ IDEA}
IntelliJ IDEA~\cite{JB:ITJ} es una herramienta que proporciona un entorno completo para el desarrollo de software, incluyendo herramientas para codificación, depuración, pruebas y despliegue.

\imagen{ITJ_logo}{Logotipo de IntelliJ IDEA }{0.3}

\subsubsection{Docker/Docker Desktop}
Docker~\cite{Docker} es un software de código abierto que automatiza el despliegue de aplicaciones dentro de contenedores de software, proporcionando una capa adicional de abstracción y automatización de visualización de aplicaciones en múltiples sistemas operativos.

En este proyecto se ha utilizado para el despliegue de ambas aplicaciones web y de una base de datos tanto Docker como Docker Desktop que la aplicación nativa para ordenador diseñada por Docker para Windows y Mac, siendo la forma más sencilla de ejecutar, construir, depurar y probar aplicaciones dockerizadas. 

Gracias a utilizar Docker se facilita el despliegue al usuario evitando que se preocupe por si tiene las dependencias necesarias o las versiones correctas de las herramientas utilizas para el correcto funcionamiento del proyecto.



\imagen{Docker_logo}{Logotipo de Docker}{0.5}

\subsubsection{Git/GitHub}
Git~\cite{Mc:Git} es un sistema de control de versiones distribuido. Estos repositorios locales plenamente funcionales permiten trabajar sin conexión o de forma remota con facilidad. Los desarrolladores confirman su trabajo localmente y, a continuación, sincronizan la copia del repositorio con la del servidor.

\imagen{Git_logo}{Logotipo de Git }{0.3}

GitHub~\cite{Github} es una plataforma basada en la nube donde puedes almacenar, compartir y trabajar junto con otros usuarios para realizar proyectos. 
Los proyectos son almacenados en repositorios, permitiendo así:
\begin{itemize}
    \item Presentar o compartir el trabajo.
    \item Seguir y administrar los cambios en el código a lo largo del tiempo.
    \item Dejar que otros usuarios revisen el código y realicen sugerencias para mejorarlo.
\end{itemize}


\imagen{Github_logo}{Logotipo de Github }{0.3}


\subsubsection{Overleaf/LaTeX}
Overleaf~\cite{Overleaf} es un editor colaborativo basado en la nube que se utiliza para escribir, editar y publicar documentos en LaTeX.

\imagen{Over_logo}{Logotipo de Overleaf}{0.3}

\subsubsection{Draw.io}
Draw.io es una herramienta gratuita y de código abierto para la creación de diagramas. Se utiliza ampliamente para diseñar diagramas de flujo, esquemas de red, diagramas UML, diagramas de arquitectura de software, organigramas y cualquier tipo de representación visual estructurada.

\imagen{Draw_logo}{Logotipo de Draw.io}{0.3}

\subsection{Backend}
\subsubsection{Python}
Python~\cite{Python} es un lenguaje de programación de alto nivel, conocido por su legibilidad y versatilidad. Se utiliza en una amplia gama de aplicaciones, desde desarrollo web hasta análisis de datos.


\imagen{Python_logo}{Logotipo de Python}{0.3}

Bibliotecas y frameworks de Python utilizados:

\begin{itemize}
    \item {Flask}: Microframework web para Python que permite desarrollar aplicaciones web de forma sencilla y modular
    \item {SQLAlchemy}: facilita la interacción con bases de datos SQL mediante objetos Python.
    \item {Pandas}: permite usar estructuras eficientes como DataFrame, ideal para la lectura, manipulación y análisis de archivos CSV.
    \item {Werkzeug}: permite usar herramientas para construir aplicaciones compatibles con WSGI, un estándar para la comunicación entre servidores web y aplicaciones Python.
    \item {Requests}: permite realizar peticiones HTTP desde el backend.
    \item {Pytz}: permite gestionar zonas horarias de forma precisa
    \item {Csv}: permite la lectura y escritura de archivos CSV.
    \item {Os}: permite interactuar con el sistema operativo (crear carpetas, comprobar rutas, acceder a archivos...).
    \item {Datetime}: permite manejar fechas y horas, como la marca temporal de los paquetes.
    \item {Traceback}: permite imprimir mensajes de error con más detalles.
    \item {Random}: permite generar datos aleatorios.
    \item {Functools}: permite funciones de orden superior y operaciones sobre objetos invocables.
    \item {mlxtend}: permite complementar otras herramientas añadiendo funcionalidades útiles para campos como la minería de datos, en este caso con la generación de reglas de asociación.  
\end{itemize}


\subsection{Frontend}

\subsubsection{HTML5/JavaScript}
HTML~\cite{Html} es el lenguaje estándar para la creación de páginas web. Define la estructura básica de los contenidos de un sitio web mediante etiquetas, que permiten organizar elementos como textos, enlaces, imágenes, tablas, formularios y otros componentes visuales.

\imagen{HTML_logo}{Logotipo de HTML 5}{0.3}

JavaScript ~\cite{JavaScript} es un lenguaje de programación interpretado, orientado a objetos y basado en prototipos, diseñado originalmente para ejecutarse en navegadores web. Su propósito principal es añadir interactividad y dinamismo a las páginas web, permitiendo modificar el contenido, responder a eventos del usuario, validar formularios, actualizar datos en tiempo real sin recargar la página.

Librerías y frameworks de Html utilizados:
\begin{itemize}
    \item {Bootstrap 5}: es un framework de CSS y JS que facilita la creación de diseños responsivos y modernos con componentes predefinidos (botones, formularios y tablas).
    \item {Chart.js}: es una biblioteca de JavaScript para generar gráficos interactivos.
    \item {Pandas}: permite usar estructuras eficientes como DataFrame, ideal para la lectura, manipulación y análisis de archivos CSV.
    \item {Jinja2}: es un motor de plantillas utilizado por Flask para insertar lógica (bucles, condiciones) y datos dinámicos en los archivos HTML.
\end{itemize}


\capitulo{5}{Aspectos relevantes del desarrollo del proyecto}

Este apartado pretende recoger los aspectos más interesantes del desarrollo del proyecto, comentados por los autores del mismo.
Debe incluir desde la exposición del ciclo de vida utilizado, hasta los detalles de mayor relevancia de las fases de análisis, diseño e implementación.
Se busca que no sea una mera operación de copiar y pegar diagramas y extractos del código fuente, sino que realmente se justifiquen los caminos de solución que se han tomado, especialmente aquellos que no sean triviales.
Puede ser el lugar más adecuado para documentar los aspectos más interesantes del diseño y de la implementación, con un mayor hincapié en aspectos tales como el tipo de arquitectura elegido, los índices de las tablas de la base de datos, normalización y desnormalización, distribución en ficheros3, reglas de negocio dentro de las bases de datos (EDVHV GH GDWRV DFWLYDV), aspectos de desarrollo relacionados con el WWW...
Este apartado, debe convertirse en el resumen de la experiencia práctica del proyecto, y por sí mismo justifica que la memoria se convierta en un documento útil, fuente de referencia para los autores, los tutores y futuros alumnos.

\capitulo{6}{Trabajos relacionados}

Actualmente el Internet de las Cosas (IoT) es un paradigma tecnológico en constante expansión que permite la conexión de dispositivos físicos a través de redes inalámbricas para la recogida, transmisión y análisis de datos en tiempo real.

Por eso mismo garantizar que la comunicación establecida en esa red es segura y robusta es un trabajo altamente solicitado y con propuestas muy diversas para llevarlo acabo.

En este apartado voy a mencionar tres trabajos relacionados con este campo que considero importante comentarlos:


\section{Desarrollo de un prototipo de modulo de análisis de
trafico basado en wireshark para detección de ataques de denegación usando inspección de tramas lorawan que provea una capa de integración api rest}

Este proyecto \cite{Rel1} se centra en la creación de un sniffer capaz de capturar y analizar paquetes LoRaWAN con el objetivo de detectar ataques de denegación de servicio (DoS/DDoS). A través del uso de herramientas como Wireshark, Pyshark y Tshark, así como del desarrollo de una API REST mediante Django, el autor consigue construir un sistema que no solo captura y decodifica paquetes LoRaWAN, sino que también los almacena en una base de datos accesible para futuras evaluaciones de seguridad.

 La propuesta de realizar una inspección de tramas en el entorno de red, es sin duda una técnica que podría aplicarse como capa adicional en mi sistema para mejorar las capacidades de diagnóstico y detección temprana de amenazas dentro del ecosistema LoRaWAN.
 


\section{Trend Micro Finds LoRaWAN Security Lacking, Develops LoRaPWN Python Utility}

Este proyecto \cite{Rel2} desarrolló LoRaPWN, una herramienta en Python diseñada para funcionar con radios definidas por software (SDR) compatibles con GNU Radio. LoRaPWN permite capturar, descifrar y generar paquetes LoRaWAN (versiones 1.0 y 1.1), romper (brute-force) el procedimiento OTAA, descifrar cargas útiles y los campos MIC, así como simular ataques DoS mediante replays y manipulaciones de paquetes.

Este proyecto esta especialmente relacionado con el mio ya que en mi proyecto se utiliza el procedimiento OTAA para unir el dispositivo IoT con TTSS, seria interesante ver como respondería mi proyecto a un intento de ataque a través del método desarrollado por este articulo, y desarrollar posibles defensas ante este tipo de ataques.



\section{Detecting IoT device compromise using Python}

El objetivo principal de este proyecto \cite{Rel3} es desarrollar una herramienta en Python capaz de analizar tráfico de red IoT a través de archivos .pcap, con el fin de detectar y modelar flujos de comunicación entre dispositivos. Esta herramienta permite extraer información útil sobre el comportamiento de los dispositivos conectados a una red local, facilitando tareas de análisis de seguridad, identificación de patrones anómalos y caracterización de dispositivos IoT.

 Esta solución permite identificar posibles vectores de ataque, caracterizar el comportamiento normal y anómalo, y servir como base para futuras estrategias de detección y mitigación de amenazas en entornos IoT que podrían aplicarse al proyecto que he desarrollado mejorando así su seguridad. 

 

\capitulo{7}{Conclusiones y Líneas de trabajo futuras}


\section{Conclusiones}
tras finalizar este proyecto y observar todo lo realizado en el mismo, puedo concluir que se han cumplido los dos objetivos generales que se tenían al comienzo del proyecto.

Se ha logrado configurar y crear una red LoRaWAN capaz de mantener un flujo real de datos y de conectar dicha red con una aplicación web capaz de recibir, analizar y almacenar esos datos para su posterior uso.

Ademas se ha logrado desarrollar una aplicación web completa, intuitiva y funcional que ayude a los usuarios no técnicos a interactuar con datos recopilados de dispositivos IoT e introducirles de forma sutil el como es el proceso de análisis y clasificación de datos con el objetivo de detectar ataques o brechas de seguridad en la red. 

Aun así creo que esa parte es la que mas margen de mejora tiene con respecto a lo desarrollado en el proyecto, ya que los métodos referentes a mejorar la seguridad de los entornos IoT es un campo de la informática que actualmente esta en auge y queda mucho por mejorar. Desde la generación de conjuntos de datos de calidad y variados para ser usados en la búsqueda de patrones o técnicas de aprendizaje automático hasta mejorar la robustez de los medios de comunicación utilizados para transmitir los paquetes de datos.

Una vez terminado el proyecto puedo afirmar que este me ha ayudado a darme cuenta de varias cosas:

La importancia de la base teórica y practica adquirida durante todo el transcurso del grado de ingeniería informática, ya que muchas cosas realizadas durante el proyecto son temas que ya había tratado con anterioridad en el grado.

La ventaja que ofrece el tener claro desde el inicio del proyecto los objetivos que deseas cumplir y las tecnologias que tienes pensdado utilizar para cumplirlos.

Por ultimo agradecer al desarrollo de este proyecto el haber aumentado mis habilidades tanto en el ámbito técnico a la hora de desarrollar software y trabajar con hardware como son los dispositivos LoRaWAN como en el ámbito organizativo a la hora de gestionar el desarrollo de un proyecto de esta magnitud y documentar todo lo relacionado con su desarrollo.


\section{Líneas de trabajo futuras}

Como ya había mencionado anteriormente y aun reconociendo los logros alcanzados con este proyecto, puedo asegurar que aun hay muchos aspectos del mismo que se pueden mejorar o nuevas funcionalidades que no se habían planteado que podría ser interesante tenerlas en cuenta:

\begin{itemize}
    \item \textbf{Mejora del dataset}:
    Ya he mencionado anteriormente los principales problemas referentes al dataset utilizado en el proyecto, por eso mismo un trabajo a futuro podría ser el mejorar la calidad del dataset generado haciéndolo mas amplio y completo. Facilitando así la tarea referente a analizar correctamente los datos recibidos.
    
    \item \textbf{Exploración de diferentes métodos de clasificación}:
    En este proyecto se ha optado por utilizar un método de clasificación que no tiene por que ser el mas adecuado para este trabajo, por ello estudiar diferentes opciones y probar si son mejores que la utilizada actualmente, mejorando así los resultados de clasificación del proyecto es una linea de trabajo a futuro completamente valida.
    
    \item \textbf{Generalización en el ámbito hardware}:
    Actualmente el proyecto desarrollado esta centrado y focalizado en dar soporte al sensor MerryIoT Motion Detection MS10, por lo tanto a futuro se podría plantear la idea de generalizar el proyecto para que sea capaz de soportar el flujo de datos de cualquier tipo de sensor IoT.
    
    \item \textbf{Análisis del flujo de datos recibido}:
    Este apartado se ha mencionado en trabajos relacionados y es la posibilidad de implementar al proyecto la capacidad de analizar el trafico de red IoT con el fin de detectar posibles vectores de ataque o comportamientos anómalos, mejorando así tanto la seguridad como el analizador del proyecto.
    
    \item \textbf{Creación de una interfaz referente a la red IoT}:
    Este apartado esta relacionado con el anterior y consiste en plantearse crear una interfaz que pueda usar el usuario, como hace en la segunda aplicación web pero referente a la primer aplicación encargada de analizar y recibir los datos, añadiendo funcionalidades adicionales como puede ser mostrar en tiempo real un grafo de nodos con el flujo de datos entre ellos.
\end{itemize}


\bibliographystyle{plain}
\bibliography{bibliografia}

\end{document}
